\begin{ozetenv}
    Büyük Dil Modellerinin (BDM) "Servis Olarak Makine Öğrenmesi" (MLaaS) platformları aracılığıyla yaygın olarak benimsenmesi, veri gizliliği konusunda önemli endişeleri beraberinde getirmiştir. Kullanıcılar, hassas bilgilerini bulut tabanlı modellerle paylaşma konusunda giderek daha fazla tereddüt yaşamaktadır. Bu tez, gizliliği koruyan doğal dil işlemeyi mümkün kılmak için Tam Homomorfik Şifreleme (FHE) ile YDM'lerin entegrasyonunu araştırmaktadır. FHE, hesaplamaların doğrudan şifrelenmiş veri üzerinde yapılmasına olanak tanıyarak, ne kullanıcının girdisinin ne de modelin ürettiği çıktının sunucuya açık metin olarak ifşa edilmemesini sağlar.
    
    Bu çalışma, Transformer tabanlı modellerin şifrelenmiş bir alanda çalıştırılması için gereken mimari adaptasyonları incelemektedir. Özellikle, Softmax ve GELU gibi doğrusal olmayan aktivasyon fonksiyonlarının, CKKS gibi FHE şemalarına uygun polinom yaklaşımlarıyla değiştirilmesi zorluklarını ele almaktadır. Araştırma; hesaplama yükü, bellek tüketimi ve çıkarım doğruluğu arasındaki ödünleşimleri değerlendirerek, güvenli YDM dağıtımı için optimize edilmiş bir çerçeve sunmaktadır.
\end{ozetenv}